% sections/01_introduction.tex --- SkillStrike

The rise of LLM agent \emph{skill marketplaces}---platforms such as
\clawhub (\NskillsCH{} skills, $>$2.1M installs) and \skillssh{}
(\NskillsSH{} skills)---has transformed language model assistants into
autonomous agents that compose third-party plugins to perform complex tasks.
A skill is an installable plugin that grants the agent ambient authority:
the ability to read files, execute shell commands, access credentials,
and make network requests, all on behalf of the user,
within a single shared context window.

Unlike traditional software stacks where each component runs in an
isolated process with OS-enforced privilege separation~\cite{felt2011android},
installed skills share the agent's entire execution context.
The agent's tool dispatcher has no awareness of which skills are
``in scope'' for a given task, and no enforcement mechanism constrains
cross-skill information or instruction flow.
This architectural property creates a structural security gap:
the \emph{aggregate capability} of a set of co-installed skills
may far exceed what any individual skill discloses,
and no user or platform mechanism surfaces this aggregate.

\paragraph{The core attack.}
We identify and formalize \emph{skill composition attacks},
a new class of agent supply-chain exploits.
The attack has three structural components:
\begin{enumerate}[leftmargin=*,topsep=2pt,itemsep=1pt]
  \item A \textbf{source} skill---installed for a benign purpose
    (web search, email, calendar, file reading)---serves as an
    \emph{injection channel}: its output reaches the agent's context
    and can carry attacker-controlled content.
  \item A \textbf{sink} skill---also installed for a legitimate purpose
    (development tooling, scheduler, memory manager)---possesses a
    \emph{dangerous undisclosed capability}: shell execution, credential
    exfiltration, boot hook injection, or persistent memory poisoning.
    This capability is absent from the skill's public documentation.
  \item The attacker poisons the source's output channel
    (a search result, email, API response, shared memory entry)
    to inject a disguised instruction that causes the agent to invoke
    the sink's dangerous capability, producing a real system-level effect
    without the user's knowledge or consent.
\end{enumerate}

Critically, \textbf{neither skill need be individually malicious}.
The attack surface arises from their composition in a shared context.
A user who audits each installed skill in isolation cannot detect
the risk, because no individual skill is performing anything unexpected---
the dangerous action emerges from the interaction between them.

\paragraph{Why skill composition attacks are structurally distinct from prompt injection.}
Prompt injection~\cite{greshake2023indirect,perez2022ignore} exploits the
model's inability to distinguish user instructions from data content.
A reviewer may ask: is a skill composition attack simply
\emph{``indirect prompt injection via tool output''}
(\\ie~\cite{zhan2024injecagent})?
We argue not, and state this precisely:

\begin{proposition}[Defense Separation]
\label{prop:separation}
Let $\mathcal{D}_{\rm PI}$ be any defense against indirect prompt injection
that operates by (a) filtering/labeling untrusted tokens before LLM processing,
(b) restricting which tokens can contain instructions, or
(c) sandboxing the attacker's injection channel.
Then $\mathcal{D}_{\rm PI}$ cannot prevent skill composition attacks
in any standard MCP-compatible agent, because:
(1) The composition attack payload travels through a \emph{trusted tool-output channel}
(a registered, installed skill's tool call result),
which $\mathcal{D}_{\rm PI}$ must exempt to preserve normal agent functionality;
(2) The dangerous tool call it triggers is a \emph{legitimate invocation}
of the sink skill's declared API---the LLM is performing no unexpected action;
(3) Blocking requires reasoning about \emph{cross-skill capability combination},
not about token-level trust---requiring a defense at the ecosystem layer,
not the input-processing layer.
\end{proposition}

\noindent
This differs from tool-output injection~\cite{zhan2024injecagent} in a critical way:
prior tool-output injection works treat the \emph{single tool's behavior} as the attack surface.
Skill composition attacks treat the \emph{interaction between two or more tools}
as the attack surface---a surface that does not exist in any single skill's schema
and cannot be detected by per-tool analysis.
The attack class is therefore not a special case of prompt injection;
it requires a distinct remediation model (capability isolation, \S\ref{sec:discussion}).

\paragraph{Contributions.}
We make four contributions:

\begin{enumerate}[leftmargin=*, label=\textbf{C\arabic*.},topsep=2pt]

\item \textbf{Attack formalization and taxonomy (\S\ref{sec:background}, \S\ref{sec:static}).}
We define the \emph{Source-Sink Composition Attack} model,
introduce a taxonomy of five sink capability families and five source injection
channels applicable across skill marketplaces,
and prove that the attack surface is $O(|S_{\rm src}| \times |S_{\rm sink}|)$
in the marketplace size.

\item \textbf{Static attack surface discovery (\S\ref{sec:static}).}
\lscps (Layered Structured Capability Pattern Search) combines
\NpatternsLSCPS{}+ regex patterns organized into 14 capability families
to identify dangerous source-sink pairs at publish time.
Applied to \NskillsCH{} \clawhub skills, \lscps identifies
\NsinkCH{} dangerous sink skills (K1--K5, 25.5\% of corpus)
and \NsrcCH{} source skills (C1--C5, 83\%),
yielding a theoretical attack surface of $\NsrcCH \times \NsinkCH \approx 71.9\mathrm{M}$ skill combinations.

\item \textbf{LLM semantic verification (\S\ref{sec:llm}).}
Applied to the \skillssh{} corpus (\NskillsSH{} skills) as a second validation corpus,
a two-stage Codex chain-of-thought audit confirms \NcodexConfirm{} exploitable paths
from \NselfExfil{} LSCPS candidates---a
$>$21$\times$ precision improvement over static analysis alone.

\item \textbf{Empirical validation on two live platforms (\S\ref{sec:eval_ch}, \S\ref{sec:eval_ss}).}
We execute \sysname attacks on two distinct, independent platforms
using real agent infrastructure (OpenClaw and Codex).
On \clawhub: \NpairsCHshell/\NpairsCHshell{} shell-execution pairs succeed
(\PairsCHshellRate), and \NpairsCHConeok/\NpairsCHCone{} LLM-context injection
pairs succeed (\PairsCHConeRate).
On \skillssh: \NpairsSSsuccess/\NpairsSSlaunched{} pairs succeed
(\PairsSSrate) across all five sink families,
with zero install-time safety gates triggered.

\end{enumerate}

\paragraph{Ethical considerations.}
All experiments ran in fully isolated sandbox environments with no impact on
real users, external services, or production skill marketplaces.
Skill identifiers used for reproducibility are real marketplace names;
no weaponized payloads are released.
We have notified both ClawHub and the skills.sh maintainers.

\paragraph{Paper organization.}
\S\ref{sec:background} defines the threat model.
\S\ref{sec:static} describes static discovery.
\S\ref{sec:llm} presents LLM verification.
\S\ref{sec:eval_ch} and \S\ref{sec:eval_ss} report empirical results.
\S\ref{sec:discussion} discusses defenses and limitations.
\S\ref{sec:related} situates our work.
